\section{Verification, remaining formal work, and excluded claims}
\label{sec:validation}

\subsection{What the executable artifacts establish}

The exact reference programs in \texttt{code/} use Python's standard
library and integer or rational arithmetic. They test the finite identities
and the algorithms used in the human arguments. Their results are recorded
in \texttt{audit/}. They are deliberately not called a proof checker for
\(\ipra\): no complete stream of first-order derivation steps is generated
and checked by these programs.

\begin{itemize}
\item The hypercube program exhaustively checks all 65,814 families in
dimensions zero through four. It verifies the edge and vertex inequalities,
the binomial thresholds, and the subcube equality count \(3^d\) for
nonempty edge-equality sets. Selected compression traces check the
potential decrease and the exceptional final configurations.
\item The support-function program checks 104 exact Laurent-polynomial
examples of the area, Wigner area, and even/odd deficit identities, and
200 exact instances of the reciprocal telescopes used in the stability
bounds. In particular it checks the one-pass Wigner area normalization.
\item The Brunn--Minkowski reference implementation checks 29 deterministic
cases, including randomized two- and three-dimensional box families with
a recorded seed. Across these examples it checks 195 tree nodes, 112
leaves, 83 volume-matching cuts, and 166 strict decreases of the child
recursion measure. It independently computes the actual Minkowski-union
volume on a common grid and verifies containment and disjointness of the
leaf sum boxes.
\item The triangle program checks the two radical-free targets, the
Heron expansion, the product-of-sides derivation, and the distinction
between nondegenerate and degenerate equality cases. Its complete counts
are recorded in its JSON output.
\end{itemize}

For illustration, one box example has \(V(X)=V(Y)=3\) and \(q=1\).
Matching a volume fraction \(1/3\) requires the rational cut \(2/3\)
in the other input. The three independently checked quantities are
\[
 \underbrace{12}_{V(X)(1+q)^2}
 \ \leq\
 \underbrace{25/2}_{\text{sum of disjoint leaf sum-box volumes}}
 \ \leq\
 \underbrace{15}_{V(X+Y)}.
\]
This demonstrates the actual cut construction, rather than only testing
the final numerical inequality.

The source was compiled to PDF, references and mathematical displays were
checked for layout problems, and representative pages were rendered for
visual inspection. The archive includes build instructions and file
hashes. The original planar source is retained separately with its source
commit and blob identifier.

\subsection{What remains for complete Lean formalization}

The task requested at least an outline, and the supplied files meet that
level: they provide definitions, specifications, internal targets, the
actual ProveIt interfaces, and a module-by-module plan. The remaining work
is substantial but identified. It includes the verified natural-number
codecs, the internal arithmetic library for the finite proofs, compilation
of the fixed algorithms into \texttt{PR} symbols, all internal
quantifier-free induction invariants, and the construction of the closed
\texttt{Prf PRAAx} derivations. The presentation transfer theorems then
require their explicit moduli and representation laws to be carried
through each theorem, rather than erased when moving to geometric notation.

No completed internal derivation is claimed merely because the outline
uses no \texttt{sorry} or added axioms. A definition of a proposition is
not a proof of it. No Lean toolchain was available for compiling these
outlines in this session. The archive's build and test successes refer
to the PDF and the finite Python diagnostics, not to a Lean build.

\subsection{Scope that should not be silently enlarged}

\begin{enumerate}
\item \textbf{All real curves or sets.} Theorems about supplied primitive
recursive names and moduli do not quantify over every abstract classical
curve or measurable set. The compatibility of enclosed area, boundary
perimeter, and the approximation data is part of the presentation.
\item \textbf{Spherical rigidity.} Cap equalities and the finite cap-union
inequality do not prove the spherical theorem for every simple closed
curve. That general theorem requires a separate effective proof.
\item \textbf{Higher-dimensional equality.} Balls attain the sharp
constant. A classification of arbitrary extremal sets, including possible
null additions under particular boundary conventions, is not derived.
\item \textbf{Hausdorff and Minkowski boundary measures.} The report uses
an explicit one-sided outer-volume modulus. General identification with
Hausdorff measure, or the existence and effectiveness of a liminf
boundary measure, is not built into that definition.
\item \textbf{Sobolev spaces.} The finite layered theorem is proved in
full. The analytic extension has an explicit coarea/norm approximation
contract. A general effective construction of that contract for every
\(W^{1,1}\) element remains separate.
\item \textbf{Hadamard manifolds and arbitrary metric spaces.} No
conversion of their general isoperimetric problems to \(\ipra\) is
asserted here. The scope decision does not depend on calling any
contemporary classical conjecture solved or unsolved; even a classical
proof would require its own constructive and arithmetic analysis.
\end{enumerate}

The new work is the explicit arithmetic packaging, proof details,
quantitative presentation contracts, and expanded coverage. The classical
inequalities themselves are established results. No claim of priority
for their underlying geometric or combinatorial ideas is intended.

